\section*{Tóm tắt}
\thispagestyle{empty}

Đồ án tốt nghiệp với đề tài “Xây dựng Hệ thống Quản lý Thư viện Trực tuyến” được thực hiện nhằm đáp ứng nhu cầu chuyển đổi số trong quản lý và khai thác tài nguyên thư viện. Trong bối cảnh các thư viện truyền thống đang từng bước chuyển mình sang môi trường số, việc xây dựng một hệ thống quản lý hiện đại, hiệu quả và hỗ trợ tốt hơn cho người dùng là yêu cầu cấp thiết. Đồ án tập trung nghiên cứu, thiết kế và triển khai một hệ thống thư viện dưới dạng Ứng dụng Web hướng tới đối tượng là các thư viện quy mô vừa và nhỏ như thư viện trường học, thư viện khoa hoặc các trung tâm học liệu.\\

Hệ thống được xây dựng nhằm quản lý đa dạng các loại tài nguyên thư viện, bao gồm sách in, tạp chí và các tài nguyên số như E-book và tài liệu PDF. Bên cạnh việc hỗ trợ các chức năng quản lý hành chính cơ bản, hệ thống hướng tới việc nâng cao trải nghiệm người dùng thông qua việc tích hợp các công nghệ hiện đại, đặc biệt là Trí tuệ nhân tạo (AI).\\

Trong hệ thống, các vai trò người dùng chính bao gồm: độc giả, thủ thư và quản trị viên. Mỗi vai trò sở hữu các yêu cầu chức năng và quy trình sử dụng riêng biệt, được hệ thống hỗ trợ một cách phù hợp:\\

Đối với vai trò độc giả, hệ thống cung cấp các chức năng tìm kiếm và tra cứu tài liệu linh hoạt thông qua từ khóa hoặc công cụ tìm kiếm thông minh. Ngoài ra, hệ thống tích hợp cơ chế gợi ý tài liệu dựa trên lịch sử đọc và hành vi người dùng, giúp độc giả nhanh chóng tiếp cận các tài liệu phù hợp với nhu cầu học tập và nghiên cứu. Độc giả có thể quản lý lịch sử mượn, theo dõi trạng thái mượn-trả và tương tác với hệ thống thông qua chatbot để được hỗ trợ kịp thời.\\

Đối với vai trò thủ thư, hệ thống hỗ trợ quản lý danh mục tài nguyên, cập nhật dữ liệu tả tài liệu và xử lý các thao tác mượn, trả sách. Hơn nữa, thủ thư có thể theo dõi tình trạng khả dụng của tài nguyên, kiểm soát dữ liệu và đảm bảo hoạt động thư viện diễn ra hiệu quả. Việc triển khai hệ thống giúp giảm thiểu các quy trình thủ công, từ đó tăng năng suất và giảm sai sót trong quản lý.\\

Đối với vai trò quản trị viên, hệ thống cung cấp quyền quản lý toàn bộ hoạt động hệ thống, bao gồm quản lý tài khoản người dùng, cấu hình hệ thống và giám sát dữ liệu. Ngoài ra, quản trị viên có thể xem các báo cáo thống kê liên quan đến hoạt động mượn-trả để hỗ trợ đưa ra quyết định quản lý và phát triển tài nguyên thư viện hợp lý.\\

Điểm nổi bật của hệ thống là việc tích hợp các tính năng Trí tuệ nhân tạo như gợi ý tài liệu, tìm kiếm bằng ngôn ngữ tự nhiên và chatbot trợ lý tương tác. Những tính năng này góp phần nâng cao đáng kể sự tương tác giữa người dùng và hệ thống, đem lại trải nghiệm vượt trội so với các hệ thống thư viện truyền thống.\\

Về mặt kiến trúc, hệ thống được thiết kế theo mô hình ứng dụng web với các thành phần chính: giao diện người dùng, logic xử lý nghiệp vụ phía máy chủ và cơ sở dữ liệu lưu trữ. Cách tổ chức này đảm bảo tính linh hoạt, khả năng mở rộng và dễ bảo trì cho các nâng cấp trong tương lai.\\

Thông qua quá trình phát triển Hệ thống Quản lý Thư viện Trực tuyến, đồ án hướng tới mục tiêu tối ưu hóa quy trình quản lý tài nguyên, nâng cao hiệu quả vận hành và mang lại trải nghiệm truy cập, phát hiện tài liệu tốt hơn cho độc giả. Đồng thời, đồ án chứng minh tính khả thi và ứng dụng thực tiễn của các công nghệ hiện đại, đặc biệt là AI, trong việc phát triển hệ thống thông tin phục vụ môi trường giáo dục và nghiên cứu.\\

\clearpage
\pagenumbering{arabic}

